\chapter{Elliptic Curve Cryptography}

By now you have met two great families of public-key cryptography. RSA leans on the
hardness of factoring; Diffie--Hellman and ElGamal lean on the discrete logarithm
problem in a multiplicative group $\Zns{p}$. Both work. Both are deployed. And both
are, by modern standards, fat: a comfortable RSA key is $3072$ bits, a comfortable
finite-field Diffie--Hellman prime is the same. The reason is the same in both cases.
The best known attacks --- the number field sieve for factoring, index calculus for
discrete logs --- run in \emph{sub-exponential} time. They are far better than brute
force, so you have to pad the key length to stay ahead of them.

Elliptic-curve cryptography is the answer to a simple question: is there a group where
the discrete logarithm problem is \emph{harder} --- where the clever index-calculus
machinery does not apply, so the only attacks left are the generic ones that work in
any group? If such a group exists, you can shrink the keys dramatically and get the
same security. The answer is yes, and the group lives on an elliptic curve. A $256$-bit
elliptic-curve key matches a $3072$-bit RSA key. That factor-of-ten saving is why your
phone, your browser, and Signal almost all run on curves today.

This chapter builds the curve group from the ground up: the geometry, the algebra over
a finite field, the worked arithmetic on tiny numbers, the elliptic-curve discrete
logarithm problem (ECDLP) and why it resists the sub-exponential attacks, then ECDH,
Curve25519, and the engineering reasons ECC won.

\chapterline

\section{From a multiplicative group to a curve}
\label{ch21:motivation}

Recall the shape of Diffie--Hellman. Fix a cyclic group $\G$ of prime order $q$ with a
generator $g$. Alice picks a secret $a$ and publishes $g^a$; Bob picks $b$ and publishes
$g^b$; the shared secret is $g^{ab} = (g^a)^b = (g^b)^a$. Security rests on the discrete
log problem: given $g$ and $g^a$, recover $a$.

Nothing in that description mentions $\Zns{p}$ specifically. Diffie--Hellman needs
\emph{any} finite cyclic group in which:

\begin{itemize}
  \item the group operation is cheap to compute,
  \item exponentiation (repeated operation) is cheap via square-and-multiply, but
  \item the discrete logarithm is hard.
\end{itemize}

The multiplicative group $\Zns{p}$ satisfies the first two. The trouble is the third:
$\Zns{p}$ has too much structure. Its elements are integers, and the index-calculus
attack exploits the fact that small integers factor into small primes, building linear
relations among discrete logs and solving the whole thing in sub-exponential time.

\begin{intuitionbox}
Index calculus is a ``the numbers know each other'' attack. In $\Zns{p}$ the elements
are ordinary integers $1,2,3,\dots$, and $6 = 2\cdot 3$ links the logarithm of $6$ to
those of $2$ and $3$. Collect enough such relations and you solve a linear system for
all the small logs, then bootstrap to any target. Points on an elliptic curve have no
analogous notion of ``factoring into smaller points.'' There is no smooth subset to
build relations from. The attack has nothing to grab onto.
\end{intuitionbox}

So we go hunting for a different group --- one whose elements are not integers and
admit no factoring story, but that still has a clean, computable group law. Elliptic
curves provide exactly that. The elements are points $(x,y)$ on a curve, and the group
operation is a geometric construction: draw a line, find where it hits the curve.

\begin{historybox}
Elliptic curves are old mathematics, studied since the nineteenth century in connection
with elliptic integrals (which compute arc lengths of ellipses --- hence the name; the
curves themselves are not ellipses). Their cryptographic use is young. In $1985$ Neal
Koblitz (University of Washington) and Victor Miller (IBM) independently proposed using
the elliptic-curve group for discrete-log cryptography. For two decades ECC was treated
as exotic and patent-encumbered; RSA dominated. The combination of expiring patents,
the rise of mobile devices that cannot spare the cycles for $3072$-bit RSA, and the
$2013$ disclosures that pushed everyone toward modern, auditable primitives finally
tipped the balance. Today new protocols default to curves.
\end{historybox}

\section{The curve and its points}
\label{ch21:curve}

Over a field $K$ whose characteristic is not $2$ or $3$ (so we can divide by $2$ and
$3$), an elliptic curve in \emph{short Weierstrass form} is the set of solutions to
\begin{equation}
\label{ch21:weierstrass}
y^2 = x^3 + ax + b, \qquad a,b \in K,
\end{equation}
together with one extra point ``at infinity,'' written $\mathcal{O}$, which we will
explain shortly. We require the cubic on the right to have no repeated roots, which is
equivalent to the \emph{nonsingularity} condition
\begin{equation}
\label{ch21:discriminant}
4a^3 + 27 b^2 \neq 0.
\end{equation}
If the discriminant vanishes the curve develops a cusp or self-crossing, the group law
breaks, and the cryptography collapses. Always check it.

\begin{figure}[H]
\centering
\begin{tikzpicture}[scale=1.05]
  \begin{axis}[
      axis lines=middle, xlabel={$x$}, ylabel={$y$},
      xmin=-2.6, xmax=3.4, ymin=-4.2, ymax=4.2,
      xtick=\empty, ytick=\empty, samples=200, domain=-1.7693:3.4,
      width=8.5cm, height=6.5cm, clip=false]
    \addplot[sourceblue, thick] {sqrt(x^3 - 2*x + 2)};
    \addplot[sourceblue, thick] {-sqrt(x^3 - 2*x + 2)};
    \node[sourceblue] at (axis cs:2.6,3.6) {$y^2=x^3-2x+2$};
  \end{axis}
\end{tikzpicture}
\caption{A real elliptic curve. It is symmetric about the $x$-axis: if $(x,y)$ is on
the curve so is $(x,-y)$, because only $y^2$ appears. This single fold of symmetry is
what gives every point an inverse.}
\label{ch21:realcurve}
\end{figure}

Over the real numbers the curve is a smooth one-dimensional shape, sometimes in one
piece and sometimes in two (an oval plus an unbounded branch). It is the symmetry
across the $x$-axis that matters: replacing $y$ by $-y$ leaves the equation untouched,
so the curve is its own mirror image. Hold onto that.

\subsection{The point at infinity}

The extra point $\mathcal{O}$ is not a typographic convenience; it is a genuine point
that makes the group law total. Picture it as sitting infinitely far up (and down) the
$y$-axis, the place where all vertical lines meet. Every vertical line ``passes through
$\mathcal{O}$.'' It will be the identity element of our group --- the analog of $1$ in
$\Zns{p}$, or of the all-zeros bit in $\xor$.

\begin{focusbox}
A clean way to make $\mathcal{O}$ rigorous is \emph{projective coordinates}. Write
points as triples $[X:Y:Z]$ with $(X,Y,Z)\ne(0,0,0)$, where $[X:Y:Z]$ and
$[\lambda X:\lambda Y:\lambda Z]$ are the same point for any $\lambda\ne 0$. The affine
point $(x,y)$ is $[x:y:1]$, and the curve becomes $Y^2 Z = X^3 + aXZ^2 + bZ^3$. Setting
$Z=0$ forces $X^3=0$, so $X=0$, leaving the single point $[0:1:0]$. That is
$\mathcal{O}$. You can carry on with the affine picture and just remember ``vertical
lines hit $\mathcal{O}$''; the projective view is there if you want the bookkeeping to
be airtight.
\end{focusbox}

\section{The group law: geometry first}
\label{ch21:grouplaw-geom}

Here is the rule that turns the set of points into a group. It is one rule, stated
geometrically, and everything else follows.

\begin{focusbox}
\textbf{Chord-and-tangent addition.} To add two points $P$ and $Q$ on the curve:
draw the straight line through $P$ and $Q$. Because the curve is a cubic, that line
meets it in exactly three points (counted with multiplicity); call the third one $R'$.
Now reflect $R'$ across the $x$-axis to get $R$. Define $P + Q = R$.
\end{focusbox}

\begin{figure}[H]
\centering
\begin{tikzpicture}[scale=1.0]
  \begin{axis}[
      axis lines=middle, xlabel={$x$}, ylabel={$y$},
      xmin=-2.7, xmax=4.2, ymin=-5.0, ymax=5.0,
      xtick=\empty, ytick=\empty, samples=200, domain=-1.7693:4.0,
      width=9.5cm, height=7.2cm, clip=false]
    \addplot[sourceblue, thick] {sqrt(x^3 - 2*x + 2)};
    \addplot[sourceblue, thick] {-sqrt(x^3 - 2*x + 2)};
    \addplot[trapred, thick, domain=-2.0:3.6] {1.0 + (1.5+x)*( (1.27-1.0)/(0.3+1.5) )};
    \node[targetgreen] at (axis cs:-1.5,1.0) {\textbullet};
    \node[targetgreen,anchor=south east] at (axis cs:-1.5,1.0) {$P$};
    \node[targetgreen] at (axis cs:0.3,1.27) {\textbullet};
    \node[targetgreen,anchor=south] at (axis cs:0.3,1.4) {$Q$};
    \node[structurepurple] at (axis cs:2.55,3.02) {\textbullet};
    \node[structurepurple,anchor=west] at (axis cs:2.62,3.02) {$R'$};
    \node[evered] at (axis cs:2.55,-3.02) {\textbullet};
    \node[evered,anchor=west] at (axis cs:2.62,-3.02) {$R=P+Q$};
    \draw[evered,dashed] (axis cs:2.55,3.02) -- (axis cs:2.55,-3.02);
  \end{axis}
\end{tikzpicture}
\caption{Adding two points. The chord through $P$ and $Q$ meets the curve at a third
point $R'$; reflecting across the $x$-axis gives $R = P+Q$.}
\label{ch21:addition-fig}
\end{figure}

Three special cases complete the definition.

\begin{itemize}
  \item \textbf{Adding $\mathcal{O}$.} For any point $P$, $P + \mathcal{O} = P$. The
        line through $P$ and ``infinity'' is the vertical line at $P$; it meets the
        curve again at the reflection $(x_P,-y_P)$, and reflecting \emph{that} back
        returns $P$. So $\mathcal{O}$ is the identity.
  \item \textbf{Inverses.} The reflection $-P = (x_P, -y_P)$ is the inverse of $P$. The
        line through $P$ and $-P$ is vertical, so it ``meets the curve at $\mathcal{O}$,''
        and reflecting $\mathcal{O}$ gives $\mathcal{O}$. Hence $P + (-P) = \mathcal{O}$.
  \item \textbf{Doubling.} To compute $P + P$ when $P=Q$, use the \emph{tangent} line at
        $P$ (the limiting position of the chord as $Q \to P$). It meets the curve at one
        other point $R'$; reflect to get $2P$.
\end{itemize}

\begin{figure}[H]
\centering
\begin{tikzpicture}[scale=1.0]
  \begin{axis}[
      axis lines=middle, xlabel={$x$}, ylabel={$y$},
      xmin=-2.7, xmax=4.2, ymin=-5.0, ymax=5.0,
      xtick=\empty, ytick=\empty, samples=200, domain=-1.7693:4.0,
      width=9.5cm, height=7.2cm, clip=false]
    \addplot[sourceblue, thick] {sqrt(x^3 - 2*x + 2)};
    \addplot[sourceblue, thick] {-sqrt(x^3 - 2*x + 2)};
    \addplot[trapred, thick, domain=-2.0:3.8] {1.414 + (x-0.0)*(-0.707)};
    \node[targetgreen] at (axis cs:0.0,1.414) {\textbullet};
    \node[targetgreen,anchor=south west] at (axis cs:0.05,1.5) {$P$};
    \node[structurepurple] at (axis cs:0.5,1.06) {\textbullet};
    \node[structurepurple,anchor=south west] at (axis cs:0.55,1.06) {$R'$};
    \node[evered] at (axis cs:0.5,-1.06) {\textbullet};
    \node[evered,anchor=west] at (axis cs:0.6,-1.2) {$2P$};
    \draw[evered,dashed] (axis cs:0.5,1.06) -- (axis cs:0.5,-1.06);
  \end{axis}
\end{tikzpicture}
\caption{Doubling a point. The tangent at $P$ meets the curve once more at $R'$;
reflecting gives $2P = P+P$.}
\label{ch21:doubling-fig}
\end{figure}

\begin{whybox}
Why does this geometric recipe form a \emph{group}? Closure and the existence of an
identity and inverses are visible from the construction. Commutativity is obvious: the
line through $P$ and $Q$ is the same as the line through $Q$ and $P$, so $P+Q=Q+P$.

The one genuinely hard fact is \textbf{associativity}: $(P+Q)+S = P+(Q+S)$. Geometrically
it is not at all obvious --- the two sides involve different chords --- yet it always
holds. The slick proof uses the theory of divisors on curves (or B\'ezout's theorem
applied to a clever pair of cubics through eight shared points). We will take
associativity as given; once you accept it, the points of the curve form an abelian
group with identity $\mathcal{O}$. That this works over \emph{any} field, including a
finite one, is the whole point.
\end{whybox}

\section{The group law: algebra}
\label{ch21:grouplaw-alg}

Geometry is the intuition; for computation we need formulas. We translate ``draw the
line, find the third intersection, reflect'' into coordinates. The arithmetic happens
in whatever field $K$ the curve lives over --- the real numbers for the pictures, but a
finite field $\F_p$ for cryptography. The formulas are identical; only the field
changes.

Let $P=(x_1,y_1)$ and $Q=(x_2,y_2)$ be points with neither equal to $\mathcal{O}$.

\subsection*{Case 1: $P \neq \pm Q$ (a genuine chord)}

The line through them has slope
\[
m = \frac{y_2 - y_1}{x_2 - x_1}.
\]
Substituting $y = m(x-x_1)+y_1$ into $y^2 = x^3+ax+b$ gives a cubic in $x$ whose three
roots are $x_1$, $x_2$, and the third intersection $x_3$. By Vieta's formula the sum of
the roots equals the coefficient of $x^2$ (which is $m^2$), so
\begin{equation}
\label{ch21:add-x}
x_3 = m^2 - x_1 - x_2, \qquad y_3 = m(x_1 - x_3) - y_1.
\end{equation}
The sign flip in $y_3$ is the reflection across the $x$-axis. Then $P+Q=(x_3,y_3)$.

\subsection*{Case 2: $P = Q$ (the tangent, ``doubling'')}

The slope of the tangent comes from implicit differentiation of the curve equation:
$2y\,dy = (3x^2+a)\,dx$, so
\[
m = \frac{3x_1^2 + a}{2y_1}.
\]
The same Vieta argument, now with a double root at $x_1$, gives
\begin{equation}
\label{ch21:double-x}
x_3 = m^2 - 2x_1, \qquad y_3 = m(x_1 - x_3) - y_1.
\end{equation}

\subsection*{Case 3: $x_1 = x_2$ but $y_1 \neq y_2$, or $y_1 = 0$}

Then $Q=-P$ (the points are vertical reflections), the line is vertical, and
$P + Q = \mathcal{O}$. Likewise if $P=Q$ and $y_1=0$ the tangent is vertical and
$2P = \mathcal{O}$.

\begin{formalbox}
Putting it together, for $P,Q \neq \mathcal{O}$ with $P \neq -Q$:
\[
m = \begin{cases}
\dfrac{y_2-y_1}{x_2-x_1} & P \neq Q,\\[1.2em]
\dfrac{3x_1^2+a}{2y_1} & P = Q,
\end{cases}
\qquad
\begin{aligned}
x_3 &= m^2 - x_1 - x_2,\\
y_3 &= m(x_1 - x_3) - y_1.
\end{aligned}
\]
(For doubling, set $x_2 = x_1$ in the $x_3$ formula, giving $m^2 - 2x_1$.) Over a finite
field $\F_p$, every operation --- the subtraction, the division, the squaring --- is
done modulo $p$, with division meaning multiplication by the modular inverse. That last
point is the only subtlety in moving from $\R$ to $\F_p$: ``divide by $x_2-x_1$'' becomes
``multiply by $(x_2-x_1)^{-1} \bmod p$,'' computed with the extended Euclidean algorithm.
\end{formalbox}

\begin{connectionbox}
The companion algorithms book covers modular inverses and the extended Euclidean
algorithm in detail. The cost of a single point addition is dominated by \emph{one}
modular inversion (to get $m$), which is the expensive part --- roughly the cost of a
modular exponentiation. This is why real implementations use \emph{projective}
coordinates: they defer all the divisions to the very end, replacing many inversions by
many cheaper multiplications. The chapter on modular exponentiation cost in that book
explains the same square-and-multiply idea we use here for scalar multiplication.
\end{connectionbox}

\section{Curves over a finite field}
\label{ch21:finitefield}

Everything above used a field only through $+,-,\times,\div$. Swap $\R$ for the finite
field $\F_p$ (the integers mod a prime $p$) and the definitions carry over verbatim. The
curve
\[
E(\F_p) = \set{(x,y) \in \F_p \times \F_p : y^2 \equiv x^3 + ax + b \pmod p} \cup \set{\mathcal{O}}
\]
is now a \emph{finite} set of points, and the chord-and-tangent law makes it a finite
abelian group. There is no picture of a smooth curve anymore --- just a scatter of dots
(Figure~\ref{ch21:scatter-fig}) --- but the algebra is unchanged.

\begin{example}[A tiny curve by hand]
\label{ch21:tinycurve}
Take $p = 13$ and the curve
\[
E : y^2 = x^3 + 3x + 8 \pmod{13}.
\]
Check nonsingularity: $4a^3 + 27b^2 = 4\cdot 27 + 27\cdot 64 = 108 + 1728 = 1836$, and
$1836 = 141\cdot 13 + 3$, so $\equiv 3 \not\equiv 0 \pmod{13}$. Good.

To list the points, for each $x \in \{0,\dots,12\}$ compute the right-hand side
$f(x)=x^3+3x+8 \bmod 13$, then ask whether it is a quadratic residue (a square) mod
$13$. The nonzero squares mod $13$ are
\[
1^2{=}1,\ 2^2{=}4,\ 3^2{=}9,\ 4^2{=}3,\ 5^2{=}12,\ 6^2{=}10,
\]
i.e. $\{1,3,4,9,10,12\}$, and each such value has two square roots $\pm y$. Tabulating:

\begin{center}
\begin{tabular}{ccccl}
\toprule
$x$ & $x^3+3x+8 \bmod 13$ & square? & $y$ values & \\
\midrule
$0$ & $8$  & no  & --- \\
$1$ & $12$ & yes & $5,8$ & ($5^2=25\equiv12$)\\
$2$ & $9$  & yes & $3,10$ \\
$3$ & $5$  & no  & --- \\
$4$ & $6$  & no  & --- \\
$5$ & $5$  & no  & --- \\
$6$ & $8$  & no  & --- \\
$7$ & $8$  & no  & --- \\
$8$ & $11$ & no  & --- \\
$9$ & $10$ & yes & $6,7$ \\
$10$& $11$ & no  & --- \\
$11$& $7$  & no  & --- \\
$12$& $4$  & yes & $2,11$ \\
\bottomrule
\end{tabular}
\end{center}

Counting: each ``yes'' row gives two points, and we add $\mathcal{O}$. There are $4$
yes-rows, so $8$ affine points plus $\mathcal{O}$, for a total of
\[
\#E(\F_{13}) = 9.
\]
\end{example}

\begin{figure}[H]
\centering
\begin{tikzpicture}[scale=0.42]
  % grid
  \draw[gray!30] (0,0) grid (12,12);
  \draw[->,gray!70] (0,0) -- (12.6,0) node[right,black]{$x$};
  \draw[->,gray!70] (0,0) -- (0,12.6) node[above,black]{$y$};
  \foreach \t in {0,3,6,9,12} {
    \node[font=\tiny,below] at (\t,0) {\t};
    \node[font=\tiny,left] at (0,\t) {\t};
  }
  % the 8 affine points of E(F_13)
  \foreach \pt in {(1,5),(1,8),(2,3),(2,10),(9,6),(9,7),(12,2),(12,11)} {
    \fill[sourceblue] \pt circle (0.28);
  }
\end{tikzpicture}
\caption{The $8$ affine points of $E:y^2=x^3+3x+8$ over $\F_{13}$ (the ninth element is
$\mathcal{O}$, off the picture). No smooth curve survives the reduction mod $p$ --- just
a scatter of dots --- yet the chord-and-tangent law still makes them a group. Note the
mirror symmetry about $y=6.5$: each point $(x,y)$ is paired with its inverse $(x,13-y)$.}
\label{ch21:scatter-fig}
\end{figure}

\begin{walkbox}
\textbf{Add two points by hand.} On the curve of Example~\ref{ch21:tinycurve}, take
$P=(1,8)$ and $Q=(2,3)$. Slope:
\[
m = \frac{3-8}{2-1} = \frac{-5}{1} = -5 \equiv 8 \pmod{13}.
\]
Then
\[
x_3 = m^2 - x_1 - x_2 = 64 - 1 - 2 = 61 \equiv 61 - 52 = 9 \pmod{13},
\]
\[
y_3 = m(x_1 - x_3) - y_1 = 8(1 - 9) - 8 = 8(-8) - 8 = -64 - 8 = -72 \equiv -72 + 78 = 6 \pmod{13}.
\]
So $P+Q = (9,6)$. Sanity check that it is on the curve: $x^3+3x+8 = 729 + 27 + 8 = 764$,
and $764 = 58\cdot 13 + 10$, so $\equiv 10$; meanwhile $y^2 = 36 \equiv 10 \pmod{13}$.
They match. The point is on the curve.
\end{walkbox}

\begin{walkbox}
\textbf{Double a point by hand.} Take $P=(1,8)$ and compute $2P$. Tangent slope:
\[
m = \frac{3x_1^2 + a}{2y_1} = \frac{3\cdot 1 + 3}{2\cdot 8} = \frac{6}{16} = \frac{6}{3}
\pmod{13},
\]
since $16 \equiv 3$. Now $6/3$: we need $3^{-1} \bmod 13$. Because $3\cdot 9 = 27 \equiv 1$,
we have $3^{-1} = 9$, so $m = 6\cdot 9 = 54 \equiv 2 \pmod{13}$. Then
\[
x_3 = m^2 - 2x_1 = 4 - 2 = 2, \qquad
y_3 = m(x_1 - x_3) - y_1 = 2(1-2) - 8 = -2 - 8 = -10 \equiv 3 \pmod{13}.
\]
So $2P = (2,3) = Q$. Interesting --- on this curve $Q = 2P$, which means
$P + Q = P + 2P = 3P$. Our earlier result $(9,6)$ is therefore $3P$.
\end{walkbox}

\begin{labbox}
Keep going from $P=(1,8)$. You found $2P=(2,3)$ and $3P=(9,6)$. Compute $4P$ by adding
$P$ to $3P$, then $5P,6P,\dots$ until you return to $\mathcal{O}$. You will find that
$P$ generates a cyclic subgroup, and its order divides $9$. (Spoiler: this $P$ has
order $9$, so $\{\mathcal{O},P,2P,\dots,8P\}$ is the \emph{entire} group --- $P$ is a
generator.) Tracing the full cycle by hand, once, fixes the whole idea in place.
\end{labbox}

\subsection{How many points? Hasse's theorem}
\label{ch21:hasse}

The group $E(\F_p)$ had $9$ points when $p=13$. Is that a coincidence, or is there a
pattern? Naively you might guess about $p$ points: for each of the $p$ values of $x$,
the right-hand side $f(x)$ is a square (two $y$'s), zero (one $y$), or a non-square (no
$y$), and ``square or not'' is roughly a coin flip, so you expect roughly $p$ affine
points plus $\mathcal{O}$. That guess is right on average, and the deviation is tightly
controlled.

\begin{theorem}[Hasse, 1933]
\label{ch21:hasse-thm}
For an elliptic curve $E$ over $\F_p$,
\[
\#E(\F_p) = p + 1 - t, \qquad \abs{t} \le 2\sqrt{p}.
\]
The integer $t$ is called the \emph{trace of Frobenius}.
\end{theorem}

So the number of points sits in the window $[\,p+1-2\sqrt p,\ p+1+2\sqrt p\,]$, of width
$4\sqrt p$ centered at $p+1$. For $p=13$ the window is $[14 - 7.2,\ 14 + 7.2] =
[6.8, 21.2]$, and indeed $9$ lands inside. For a $256$-bit prime the group order is a
$256$-bit number known to within a relative error of about $2^{-127}$ --- essentially
$p$ itself.

\begin{whybox}
Hasse's theorem is the quantitative reason ECC ``just works.'' It guarantees the group
is big (roughly $p$ points), so the discrete-log search space is huge, and it lets a
designer compute the exact order with a polynomial-time algorithm (Schoof's algorithm
counts points without enumerating them). Knowing the exact order matters because secure
ECC wants the order to be prime, or a small cofactor times a large prime. If the order
factors into many small pieces, the Pohlig--Hellman attack (Section~\ref{ch21:ecdlp})
chops the discrete log into easy sub-problems and the curve is dead.
\end{whybox}

\section{Scalar multiplication and the ECDLP}
\label{ch21:ecdlp}

In a multiplicative group we wrote the repeated operation as exponentiation $g^a$. On a
curve the operation is written additively, so ``repeating'' means \emph{scalar
multiplication}:
\[
kP = \underbrace{P + P + \cdots + P}_{k \text{ times}}.
\]
This is the exact analog of $g^k$. And just as $g^k$ is computed by square-and-multiply,
$kP$ is computed by \emph{double-and-add}: write $k$ in binary and process the bits.

\begin{example}[Double-and-add]
\label{ch21:dbladd}
To compute $13P$, note $13 = 1101_2$. Starting from $\mathcal{O}$ and reading bits left
to right, double at each step and add $P$ where the bit is $1$:
\[
\mathcal{O} \xrightarrow{\,1\,} P \xrightarrow{\,1\,} 2P{+}P{=}3P
\xrightarrow{\,0\,} 6P \xrightarrow{\,1\,} 12P{+}P{=}13P.
\]
Four doublings and two additions, instead of twelve additions. For a $256$-bit scalar
this is about $256$ doublings and $128$ additions --- a few hundred curve operations ---
versus $2^{256}$ for the naive method.
\end{example}

\begin{definition}[Elliptic-Curve Discrete Logarithm Problem]
\label{ch21:ecdlp-def}
Let $E$ be an elliptic curve over $\F_p$ and let $P \in E(\F_p)$ be a point of large
prime order $n$. Given $P$ and $Q = kP$ for an unknown integer $k \in \{0,\dots,n-1\}$,
the \textbf{ECDLP} is to recover $k$.
\end{definition}

Computing $kP$ from $k$ is easy (double-and-add). Going backward --- recovering $k$ from
$Q$ --- is believed to be hard. That asymmetry is the entire cryptographic engine, the
exact mirror of the discrete log problem in $\Zns{p}$.

\begin{intuitionbox}
Think of scalar multiplication as a deterministic but wild walk around the curve.
Starting from $P$ and repeatedly stepping, you visit $2P, 3P, 4P, \dots$ in an order
that looks random as a function of the coordinates. Given a destination $Q$, there is no
shortcut that reads off the number of steps; you appear to have to walk. The best
generic walk is a clever cycle-finding trick (Pollard's rho) that needs about
$\sqrt{n}$ steps, not $n$ --- but $\sqrt{n}$ is still astronomically large when $n$ is
$256$ bits.
\end{intuitionbox}

\subsection{Why ECC keys are small: generic versus special attacks}
\label{ch21:keysize}

Here is the crux of the whole chapter. In \emph{any} finite cyclic group of prime order
$n$, the discrete log can be solved by generic algorithms --- baby-step giant-step or
Pollard's rho --- in time about $\sqrt{n}$. These attacks treat the group as a black box;
they work the same on $\Zns{p}$, on a curve, on anything. The $\sqrt{n}$ bound is a hard
floor: Shoup proved that no generic algorithm can do better.

The difference between $\Zns{p}$ and an elliptic curve is what happens \emph{beyond} the
generic attacks.

\begin{itemize}
  \item In $\Zns{p}$, the elements are integers, and \textbf{index calculus} exploits
        their factorization structure to solve discrete logs in \emph{sub-exponential}
        time $L_p[1/3] = \exp\!\big(c\,(\ln p)^{1/3}(\ln\ln p)^{2/3}\big)$. This is far
        faster than $\sqrt n$. To stay secure you must inflate $p$ until even
        sub-exponential time is infeasible --- hence $3072$-bit primes.
  \item On a well-chosen elliptic curve, \textbf{no index-calculus analog is known}.
        There is no factoring story for points, no smooth subset, nothing to build
        relations from. The best known attack is the generic $\sqrt n$. So $n$ only has
        to be big enough to make $\sqrt n$ infeasible.
\end{itemize}

\begin{whybox}
``Infeasible'' means roughly $2^{128}$ operations --- the same security level AES-128
targets. For the generic attack to cost $2^{128}$ we need $\sqrt n \approx 2^{128}$,
i.e. $n \approx 2^{256}$. That is why a curve over a $256$-bit field gives $128$-bit
security with a $256$-bit key.

For $\Zns{p}$ to reach $128$-bit security against sub-exponential index calculus, you
solve $L_p[1/3] \approx 2^{128}$, which lands around $p \approx 2^{3072}$. Same security,
$12\times$ the key length and far more arithmetic per operation. The table makes the
trade stark.
\end{whybox}

\begin{center}
\begin{tabular}{lccc}
\toprule
Security level & Symmetric (AES) & RSA / finite-field DH & Elliptic curve \\
\midrule
$80$  & $80$  & $1024$ & $160$ \\
$112$ & $112$ & $2048$ & $224$ \\
$128$ & $128$ & $3072$ & $256$ \\
$192$ & $192$ & $7680$ & $384$ \\
$256$ & $256$ & $15360$ & $512$ \\
\bottomrule
\end{tabular}
\end{center}

(All entries in bits; these are the long-standing NIST equivalences.) The right column
grows linearly with the security level while the RSA column explodes. That gap is why a
$521$-bit curve handles top-secret-grade traffic while the RSA equivalent needs
$15360$ bits and is painfully slow.

\begin{connectionbox}
The $\sqrt n$ figure for Pollard's rho is a \emph{birthday bound} in disguise. The
algorithm generates a pseudo-random sequence of points and waits for a collision; by the
birthday paradox a collision appears after about $\sqrt n$ steps, and a collision yields
a linear equation that reveals $k$. The companion probability book derives the
$\sqrt{\text{(space size)}}$ birthday bound in general; here the ``space'' is the group
of size $n$. Same mathematics, applied to break a discrete log.
\end{connectionbox}

\subsection{When the security collapses: pitfalls}
\label{ch21:pitfalls}

ECC is only as strong as the curve and the implementation. Several traps have bitten
real systems.

\begin{attackbox}
\textbf{Pohlig--Hellman: composite group order.} If the point's order $n$ factors as
$n = \prod n_i$ with small $n_i$, the ECDLP splits (by the Chinese remainder theorem)
into independent sub-problems, one modulo each $n_i$. Each is solvable in $\sqrt{n_i}$
time, and you reassemble $k$ from the pieces. A curve whose order is, say,
$2^3 \cdot 3 \cdot (\text{500-bit prime})$ is only as strong as that $500$-bit prime ---
fine --- but a curve whose order is a product of many medium primes is broken. \emph{Fix:}
choose the point's order to be a large prime $n$, with at most a tiny cofactor
$h = \#E/n$ (typically $h \in \{1,2,4,8\}$).
\end{attackbox}

\begin{attackbox}
\textbf{Anomalous and supersingular curves.} A handful of special curve families have an
``extra'' structure that defeats the generic floor. \emph{Anomalous} curves, where
$\#E(\F_p) = p$ exactly, admit the Smart--Satoh--Araki attack: a map to the additive
group $\F_p$ solves the discrete log in \emph{linear} time. \emph{Supersingular} curves
admit the MOV/Frey--R\"uck attack, which uses a pairing to push the ECDLP into a small
extension field $\F_{p^k}$ where index calculus applies. Standard curves are screened to
avoid both. \emph{Fix:} use a vetted curve, never one you rolled yourself.
\end{attackbox}

\begin{attackbox}
\textbf{Invalid-curve and small-subgroup attacks.} If your code receives a point from
the network and does \emph{not} verify it actually lies on the agreed curve, an attacker
can send a point on a different, weak curve (or a low-order point). The arithmetic
formulas only use $a$ (not $b$), so they happily compute on the wrong curve, and the
attacker reads off your secret modulo a small number. Repeated, this recovers the whole
key. This broke real TLS and Bluetooth implementations. \emph{Fix:} always check that a
received point satisfies the curve equation and is not $\mathcal{O}$ or low-order; or use
a curve (like X25519) whose design neutralizes the attack.
\end{attackbox}

\begin{trapbox}
The thread through every one of these is: \textbf{do not invent your own curve, and do
not skip point validation.} The mathematics of secure curve selection (point counting,
cofactor analysis, twist security, avoiding anomalous/supersingular families) is
specialist work. Use NIST P-256, Curve25519, or another standardized curve, and validate
every point you receive.
\end{trapbox}

\section{Elliptic-Curve Diffie--Hellman (ECDH)}
\label{ch21:ecdh}

With a hard ECDLP in hand, Diffie--Hellman transplants onto the curve with zero
conceptual change. Replace the cyclic group $\Zns{p}$ by the curve subgroup
$\langle P\rangle$ of prime order $n$, replace exponentiation $g^a$ by scalar
multiplication $aP$, and you have \textbf{ECDH}.

\begin{focusbox}
\textbf{The ECDH protocol.} Public parameters: a curve $E/\F_p$ and a base point $P$ of
prime order $n$.
\begin{enumerate}
  \item Alice picks a secret scalar $a \getsr \{1,\dots,n-1\}$ and sends $A = aP$.
  \item Bob picks a secret scalar $b \getsr \{1,\dots,n-1\}$ and sends $B = bP$.
  \item Alice computes $S = aB = a(bP) = abP$.
  \item Bob computes $S = bA = b(aP) = abP$.
\end{enumerate}
Both arrive at the same point $S = abP$. They derive a symmetric key by hashing $S$
(specifically its $x$-coordinate): $k = \Hash(x_S)$.
\end{focusbox}

\begin{figure}[H]
\centering
\begin{tikzpicture}[font=\small, >={Stealth[round]}]
  \node[draw, rounded corners, fill=aliceblue!12, minimum width=3.2cm, minimum height=1.0cm] (A) at (0,0) {\textbf{Alice}};
  \node[draw, rounded corners, fill=bobgreen!12, minimum width=3.2cm, minimum height=1.0cm] (B) at (8.5,0) {\textbf{Bob}};
  \node[below=2pt of A, align=center, text=aliceblue] {secret $a$\\ $A = aP$};
  \node[below=2pt of B, align=center, text=bobgreen] {secret $b$\\ $B = bP$};
  \draw[->, thick] ($(A.east)+(0,0.25)$) -- node[above]{$A = aP$} ($(B.west)+(0,0.25)$);
  \draw[<-, thick] ($(A.east)+(0,-0.25)$) -- node[below]{$B = bP$} ($(B.west)+(0,-0.25)$);
  \node[below=1.7cm of A, align=center, text=aliceblue] {compute\\ $aB = abP$};
  \node[below=1.7cm of B, align=center, text=bobgreen] {compute\\ $bA = abP$};
  \node[draw, rounded corners, fill=structurepurple!12, minimum width=4.0cm] (S) at (4.25,-3.4) {shared secret $S = abP$};
  \draw[->, dashed] ($(A.south)+(0,-1.9)$) |- (S.west);
  \draw[->, dashed] ($(B.south)+(0,-1.9)$) |- (S.east);
\end{tikzpicture}
\caption{ECDH message flow. Each party sends its public point; each multiplies the
other's point by its own secret scalar; both land on $abP$.}
\label{ch21:ecdh-fig}
\end{figure}

\begin{whybox}
Eve sees $P$, $A=aP$, and $B=bP$. To get $S = abP$ she must either recover $a$ from $A$
or $b$ from $B$ --- both ECDLP instances --- or solve the \emph{Computational
Diffie--Hellman} problem on the curve directly: produce $abP$ from $aP$ and $bP$ without
either scalar. No method is known to do this faster than solving an ECDLP, which on a
good $256$-bit curve costs about $2^{128}$ operations. So ECDH inherits exactly the
security of the underlying ECDLP. As with classic Diffie--Hellman, plain ECDH gives
\emph{no} authentication --- an active attacker can man-in-the-middle it --- so it is
always paired with signatures or a pre-shared identity (the topic of later chapters on
authenticated key exchange).
\end{whybox}

\begin{example}[ECDH on the toy curve]
\label{ch21:ecdh-toy}
Reuse $E: y^2 = x^3+3x+8$ over $\F_{13}$ with base point $P=(1,8)$ of order $9$. Say
Alice's secret is $a=2$ and Bob's is $b=3$.
\begin{itemize}
  \item Alice sends $A = 2P = (2,3)$ (computed earlier).
  \item Bob sends $B = 3P = (9,6)$.
  \item Alice computes $aB = 2\cdot(3P) = 6P = (9,7)$.
  \item Bob computes $bA = 3\cdot(2P) = 6P = (9,7)$.
\end{itemize}
Both get $S = 6P = (9,7)$, agreeing as required. (In a real system $n$ is a $256$-bit
prime, not $9$, so the secret scalars are astronomically large --- but the mechanism is
identical.)
\end{example}

\begin{implementbox}
A toy ECDH in Python, written for clarity, not safety. The arithmetic mirrors the
formulas of Section~\ref{ch21:grouplaw-alg} exactly.

\begin{lstlisting}[language=Python]
# TOY elliptic-curve Diffie-Hellman -- FOR LEARNING ONLY.
# No constant-time arithmetic, no point validation, no side-channel
# protection, tiny insecure parameters. NEVER use this for anything real.

p, a, b = 13, 3, 8                 # curve y^2 = x^3 + 3x + 8 over F_13
O = None                           # the point at infinity (identity)

def inv(z):                        # modular inverse via Fermat (p prime)
    return pow(z % p, p - 2, p)

def add(P, Q):
    if P is O: return Q
    if Q is O: return P
    x1, y1 = P; x2, y2 = Q
    if x1 == x2 and (y1 + y2) % p == 0:
        return O                   # P == -Q  ->  infinity
    if P == Q:
        m = (3 * x1 * x1 + a) * inv(2 * y1) % p     # tangent slope
    else:
        m = (y2 - y1) * inv(x2 - x1) % p            # chord slope
    x3 = (m * m - x1 - x2) % p
    y3 = (m * (x1 - x3) - y1) % p
    return (x3, y3)

def scalar_mul(k, P):              # double-and-add
    R, base = O, P
    while k:
        if k & 1: R = add(R, base)
        base = add(base, base)
        k >>= 1
    return R

P = (1, 8)                         # base point, order 9 on this curve
a_sec, b_sec = 2, 3                # Alice's and Bob's secret scalars
A = scalar_mul(a_sec, P)           # Alice's public point
B = scalar_mul(b_sec, P)           # Bob's public point
print("Alice:", scalar_mul(a_sec, B))   # a*(bP)
print("Bob:  ", scalar_mul(b_sec, A))   # b*(aP)  -- same point
\end{lstlisting}

Running this prints the same shared point for both parties. The warning is not
boilerplate: a real ECDH must use constant-time scalar multiplication (to resist the
timing attacks of a later chapter), must validate incoming points (to resist the
invalid-curve attack above), and must use a vetted curve with a cryptographic field size.
\end{implementbox}

\section{Curve25519 and X25519}
\label{ch21:curve25519}

The most widely deployed curve today is not one of the original NIST curves but
\textbf{Curve25519}, designed by Daniel J. Bernstein in 2005. It is the curve behind the
\textbf{X25519} key-exchange function used in TLS 1.3, the Signal protocol, SSH, WireGuard,
and countless others. Understanding why it displaced the NIST curves is a lesson in
practical cryptographic engineering.

Curve25519 is a \emph{Montgomery curve},
\[
E : y^2 = x^3 + 486662\,x^2 + x \quad \text{over}\quad \F_p, \qquad p = 2^{255} - 19,
\]
a slightly different normal form from short Weierstrass but with the same group law in
spirit. The prime $p = 2^{255}-19$ is chosen to be just under a power of two, which makes
reduction modulo $p$ extremely fast on real hardware. The curve order is $8n$ where $n$
is a large prime --- cofactor $8$ --- and the design swallows that cofactor cleanly.

\begin{aestheticbox}
What makes X25519 special is not raw speed (though it is fast) but how aggressively it
removes \emph{footguns}. Several historically dangerous decisions are taken out of the
implementer's hands:
\begin{itemize}
  \item \textbf{Montgomery ladder.} Scalar multiplication uses a ladder that performs the
        same operations regardless of the secret bits, making constant-time
        implementation natural --- timing side channels are designed out.
  \item \textbf{$x$-coordinate only.} X25519 works with just the $x$-coordinate of points.
        This shrinks keys to $32$ bytes and means every $32$-byte string maps to a valid
        $x$, so there is no ``is this point on the curve?'' check to forget --- the
        invalid-curve attack is structurally prevented.
  \item \textbf{Cofactor handling.} ``Clamping'' the scalar (forcing certain bits)
        guarantees the result lands in the large prime-order subgroup, neutralizing
        small-subgroup attacks.
  \item \textbf{Rigid design.} The constants are derived from clear criteria with little
        wiggle room, reducing the worry that a curve was secretly chosen to be weak ---
        a concern that dogged the NIST curves after 2013.
\end{itemize}
\end{aestheticbox}

\begin{focusbox}
\textbf{Using X25519 (the right way).} As an application programmer you never touch the
group law. The interface is two functions:
\begin{itemize}
  \item Generate a $32$-byte random secret $a$; compute the public key
        $A = \texttt{X25519}(a, \text{basepoint})$.
  \item On receiving the peer's public key $B$, compute the shared secret
        $S = \texttt{X25519}(a, B)$, then run $S$ through a key-derivation function (HKDF)
        to get your symmetric keys.
\end{itemize}
That is the whole API. The library handles clamping, the ladder, and reduction. This is
the model you should follow everywhere in cryptography: use a vetted, narrow,
hard-to-misuse interface, and never reach for the primitives directly.
\end{focusbox}

\begin{applicationbox}
\textbf{Where curves run the internet.}
\begin{itemize}
  \item \textbf{TLS 1.3} (HTTPS): the key exchange is ECDH, and X25519 is the default
        group in essentially every modern browser and server. When you load a web page,
        an elliptic-curve Diffie--Hellman almost certainly established the session key.
  \item \textbf{Signal / WhatsApp}: the X3DH handshake and the Double Ratchet both run on
        X25519, and identity signatures use Ed25519 (the signature scheme on the same
        curve, covered in the digital-signatures chapter).
  \item \textbf{SSH}: \texttt{curve25519-sha256} is the default key exchange in OpenSSH.
  \item \textbf{WireGuard} (modern VPNs): X25519 is the \emph{only} key-exchange
        primitive --- the protocol is deliberately built around a single hard-to-misuse
        curve.
  \item \textbf{Cryptocurrencies}: Bitcoin and Ethereum sign transactions with ECDSA on
        the curve \texttt{secp256k1}.
\end{itemize}
The common thread: small keys, fast operations, and side-channel-resistant designs make
curves the right fit for phones, embedded devices, and high-volume servers alike.
\end{applicationbox}

\begin{historybox}
The shift toward Curve25519 accelerated after 2013, when documents revealed that the
constant $P$ used in the NIST-standardized random generator \texttt{Dual\_EC\_DRBG}
appeared to contain a backdoor. That was a different primitive from the NIST elliptic
curves, but it poisoned trust in any standard whose constants could not be fully
explained. Curve25519's rigid, criteria-driven design --- where the constants leave
almost no room to hide a weakness --- became a major selling point, and the IETF adopted
it as a first-class option. It is now arguably the most-used curve on Earth.
\end{historybox}

\section{Why ECC dominates: the summary}
\label{ch21:dominance}

Pull the threads together. Elliptic-curve cryptography wins on four fronts at once.

\begin{enumerate}
  \item \textbf{Small keys for equal security.} $256$-bit ECC matches $3072$-bit RSA.
        Less bandwidth per handshake, less storage per key, smaller certificates.
  \item \textbf{Fast operations.} Scalar multiplication on a $256$-bit curve is cheaper
        than a $3072$-bit modular exponentiation, so handshakes are faster and servers
        carry more connections per core.
  \item \textbf{Hard problem, no shortcut.} The ECDLP resists the sub-exponential
        index-calculus attacks that force RSA and finite-field DH to inflate their keys;
        the best attack on a good curve is generic $\sqrt n$.
  \item \textbf{Misuse-resistant modern designs.} Curves like Curve25519 engineer away
        timing leaks, invalid-point attacks, and small-subgroup attacks at the level of
        the primitive, not the application.
\end{enumerate}

\begin{marginnotebox}
One caveat for the road ahead: ECC, like RSA and finite-field DH, falls to a large-scale
quantum computer. Shor's algorithm solves the ECDLP (and factoring, and finite-field
discrete log) in polynomial time on a quantum machine. The smaller key does not help ---
in fact a $256$-bit curve needs \emph{fewer} logical qubits to break than a $3072$-bit
RSA modulus. The post-quantum chapter takes up what replaces curves when that day
arrives; current practice is to run X25519 \emph{alongside} a post-quantum scheme in a
hybrid handshake, so the session is safe if either one holds.
\end{marginnotebox}

\section*{What to remember}

An elliptic curve over a finite field is a finite abelian group whose elements are
points $(x,y)$ satisfying $y^2 = x^3 + ax + b$, plus a point at infinity $\mathcal{O}$
serving as the identity. The group operation is chord-and-tangent addition, captured by
explicit formulas built from one modular inversion. Repeated addition is scalar
multiplication $kP$, computed by double-and-add; reversing it is the ECDLP. Because no
index-calculus attack is known on a well-chosen curve, the ECDLP is harder per bit than
discrete logs in $\Zns{p}$, so $256$-bit curves deliver $128$-bit security where RSA
needs $3072$ bits. ECDH transplants Diffie--Hellman onto the curve unchanged, and
X25519 packages it in a fast, hard-to-misuse interface that now secures most of the
internet's key exchanges.

\section{Exercises}
\label{ch21:exercises}

\begin{exercise}[Point counting by hand]
\label{ch21:ex-count}
Let $E: y^2 = x^3 + 2x + 3$ over $\F_7$. First verify nonsingularity via
$4a^3 + 27b^2 \pmod 7$. Then, for each $x \in \{0,\dots,6\}$, compute $f(x)=x^3+2x+3
\bmod 7$, decide whether it is a quadratic residue (the nonzero squares mod $7$ are
$\{1,2,4\}$), and list all points. How many points does $E(\F_7)$ have, and does your
count satisfy Hasse's bound $\abs{\#E - 8} \le 2\sqrt 7 \approx 5.29$?
\end{exercise}

\begin{exercise}[Adding and doubling]
\label{ch21:ex-add}
On the curve $E: y^2 = x^3 + 3x + 8$ over $\F_{13}$, take $P=(1,8)$. Using the algebraic
formulas, compute $4P$ (add $P$ to $3P=(9,6)$), then $5P,\dots$ until you reach
$\mathcal{O}$. Confirm that $9P = \mathcal{O}$, proving $P$ has order $9$. Show all
modular-inverse computations. (You should find $4P=(12,11)$ and $5P=(12,2)$.)
\end{exercise}

\begin{exercise}[A full group]
\label{ch21:ex-subgroup}
Continuing the previous exercise, write out the complete cyclic group
$\langle P\rangle = \{\mathcal{O}, P, 2P, \dots, 8P\}$ with coordinates (here $P$
generates all nine points). Verify that $4P + 5P = \mathcal{O}$ and that $-P = 8P$ by
checking that $8P = (x_P, -y_P)$.
\end{exercise}

\begin{exercise}[Toy ECDH]
\label{ch21:ex-ecdh}
Using the curve and base point of Exercise~\ref{ch21:ex-add}, let Alice's secret be
$a=3$ and Bob's be $b=4$. Compute the public points $A=3P$ and $B=4P$, then the shared
secret from each side ($3B$ and $4A$), and confirm they agree. What goes wrong with this
``key exchange'' for real use, given that $P$ has order only $9$?
\end{exercise}

\begin{exercise}[Why no eavesdropper shortcut]
\label{ch21:ex-eve}
Explain, in terms of the ECDLP and the Computational Diffie--Hellman problem, exactly
which computational problem Eve must solve to recover the ECDH shared secret from the
transcript $(P, aP, bP)$. Why does knowing $aP$ and $bP$ not let her ``multiply'' to get
$abP$ directly?
\end{exercise}

\begin{exercise}[Key-size reasoning]
\label{ch21:ex-keysize}
The best generic discrete-log attack runs in time about $\sqrt n$ where $n$ is the
subgroup order. (a) What size $n$ (in bits) gives $128$-bit security? (b) Explain why
$\Zns{p}$ needs a roughly $3072$-bit prime for the same $128$-bit security, naming the
attack responsible. (c) A junior engineer proposes a $128$-bit curve for ``$128$-bit
security.'' What is the actual security level of a curve with a $128$-bit group order,
and why?
\end{exercise}

\begin{exercise}[Pohlig--Hellman intuition]
\label{ch21:ex-pohlig}
Suppose a point $P$ has order $n = 2^4 \cdot 3 \cdot 5 \cdot 7 \cdot (\text{a }30\text{-bit prime})$.
Roughly how much work is the ECDLP in $\langle P\rangle$, and which factor dominates?
State the design rule this motivates for choosing the group order, and identify which
attack class it defends against.
\end{exercise}

\begin{exercise}[Invalid-curve attack]
\label{ch21:ex-invalid}
The addition formulas of Section~\ref{ch21:grouplaw-alg} never use the coefficient $b$.
Explain how an attacker exploits this when a server fails to check that a received point
lies on the agreed curve, and why working ``on the wrong curve'' can leak the server's
secret scalar modulo small numbers. State the two-line fix.
\end{exercise}

\begin{exercise}[Discriminant matters]
\label{ch21:ex-singular}
Show that the curve $y^2 = x^3$ (so $a=b=0$) fails the nonsingularity test, and find the
single point where it is singular. Explain geometrically why the chord-and-tangent
construction breaks there (hint: what does the curve look like near the origin?).
\end{exercise}

\begin{exercise}[Prove an inverse formula]
\label{ch21:ex-inverse}
Prove from the group law that for any non-identity point $P=(x,y)$ on
$y^2 = x^3+ax+b$, the inverse is exactly $-P = (x,-y)$. (Show that $P + (x,-y)$ falls
into the vertical-line case and therefore equals $\mathcal{O}$.) Then prove $-\mathcal{O}
= \mathcal{O}$.
\end{exercise}

\begin{exercise}[Why X25519 needs no point check]
\label{ch21:ex-x25519}
X25519 operates on only the $x$-coordinate of points and accepts any $32$-byte string as
input. Argue informally why this design choice removes the invalid-curve attack of
Exercise~\ref{ch21:ex-invalid} as a class --- that is, why there is no ``off-curve''
input an attacker can supply that the implementation must reject. (Hint: think about the
quadratic twist and how clamping confines the result to the safe subgroup.)
\end{exercise}
